\section{Abstract}\label{abstract}

This note closes the current stage of the closure-node / redox-richness
programme. The initial record, \emph{Closure-Node Registration},
proposed that redox richness can be read as where unresolved relational
difference is registered relative to an electronic closure node. In the
transition-metal seed dataset, the descriptor

\[
d_{\mathrm{balance}}=1-|N_d-5|/5
\]

organizes oxidation-state count and span better than linear or monotonic
baselines. The post-audit conclusion is narrower: this is a compact
quantitative encoding of the known half-filled-shell pattern, not a
demonstrated predictor superior to a fair quadratic baseline in the same
filling coordinate.

The subsequent actinide and radius--complement tests sharpen the
epistemic status further. Static \(f\)-balance, static \(s:d:f\)
fractions, and a quadratic radius--complement representation all fail to
establish predictive novelty once degree-of-freedom-matched baselines
and label audits are applied. A first non-circular redox-couple
potential test also fails to show predictive advantage. The programme
should therefore be framed as a \textbf{generative mechanism framework}:
it creates relational coordinates, hypotheses, and falsifiable
experimental designs, but it is not yet a validated predictor of
chemical redox behaviour.

\section{1. Position}\label{position}

The correct current claim is:

\[
\boxed{\text{The framework generates interpretable relational descriptions and falsifiable tests; it is not yet a validated predictive model.}}
\]

The descriptor and geometry developed so far are useful as
\emph{mechanism screens}. They are not yet useful as independently
validated predictors. This distinction should be explicit in the
abstract, introduction, and conclusion of any public article.

\section{2. Experiment 1: Closure-node registration
record}\label{experiment-1-closure-node-registration-record}

The first associated experiment is the Zenodo v1 record:

\begin{itemize}
\tightlist
\item
  Title: \emph{Closure-Node Registration: Redox richness as the
  signature of where unresolved relational difference is registered
  relative to a closure node --- and two regimes of one event in the d-
  and f-blocks}.
\item
  DOI: \texttt{10.5281/zenodo.20573845}.
\item
  Published: 8 June 2026.
\item
  Resource type: preprint.
\item
  Files: \texttt{CLOSURE\_NODE\_REGISTRATION.pdf} and
  \texttt{ROSI\_article\_support\_2026-06-06.zip}.
\end{itemize}

The original d-block result remains reproducible from the archived
support data: \texttt{d\_balance} obtains LOSO-RMSE \(1.2259\) for
oxidation-state count and \(1.4755\) for oxidation span, versus
\(1.6483\) and \(1.9708\) for a linear \(N_d\) baseline.

However, the fair comparison is not against linear \(N_d\). Since
\(d_{\mathrm{balance}}\) is an affine transform of \(|N_d-5|\), it
already imposes a tent-shaped peak at half filling. The fair baseline is
a quadratic in \(N_d\), or a quadratic in \(|N_d-5|\).

\section{3. Resolution of the lyo\_bench
critique}\label{resolution-of-the-lyo_bench-critique}

The Reddit critique correctly identifies four required fixes:

\begin{enumerate}
\def\labelenumi{\arabic{enumi}.}
\tightlist
\item
  compare against quadratic filling baselines, not only linear \(N_d\);
\item
  bootstrap model-comparison error, not only the \(d^5\) local mean
  contrast;
\item
  respect the leave-one-series-out design by block/cluster bootstrapping
  at the series level;
\item
  lock the oxidation-state source and model set before fitting.
\end{enumerate}

The series-block bootstrap below uses the same three blocks as the LOSO
design: \(3d\), \(4d\), and \(5d\). With only three clusters the
intervals are still coarse, but they are more honest than element-wise
resampling.

\subsection{3.1 Count models}\label{count-models}

\begin{longtable}[]{@{}lll@{}}
\toprule\noalign{}
model & LOSO\_RMSE & AICc\_in\_sample \\
\midrule\noalign{}
\endhead
\bottomrule\noalign{}
\endlastfoot
d\_balance & 1.226 & 10.743 \\
Nd linear & 1.648 & 34.641 \\
Nd quadratic & 1.291 & 11.700 \\
& Nd-5 & quadratic \\
group & 1.657 & 34.616 \\
period+group & 1.751 & 36.303 \\
\end{longtable}

\subsection{3.2 Span models}\label{span-models}

\begin{longtable}[]{@{}lll@{}}
\toprule\noalign{}
model & LOSO\_RMSE & AICc\_in\_sample \\
\midrule\noalign{}
\endhead
\bottomrule\noalign{}
\endlastfoot
d\_balance & 1.475 & 24.950 \\
Nd linear & 1.971 & 45.061 \\
Nd quadratic & 1.509 & 22.701 \\
& Nd-5 & quadratic \\
group & 1.993 & 45.056 \\
period+group & 1.994 & 46.574 \\
\end{longtable}

\subsection{3.3 Series-block bootstrap of RMSE
differences}\label{series-block-bootstrap-of-rmse-differences}

Here

\[
\Delta RMSE = RMSE(d_{\mathrm{balance}}) - RMSE(\text{baseline}).
\]

Negative values favour \texttt{d\_balance}.

\begin{longtable}[]{@{}
  >{\raggedright\arraybackslash}p{(\columnwidth - 10\tabcolsep) * \real{0.1667}}
  >{\raggedright\arraybackslash}p{(\columnwidth - 10\tabcolsep) * \real{0.1667}}
  >{\raggedright\arraybackslash}p{(\columnwidth - 10\tabcolsep) * \real{0.1667}}
  >{\raggedright\arraybackslash}p{(\columnwidth - 10\tabcolsep) * \real{0.1667}}
  >{\raggedright\arraybackslash}p{(\columnwidth - 10\tabcolsep) * \real{0.1667}}
  >{\raggedright\arraybackslash}p{(\columnwidth - 10\tabcolsep) * \real{0.1667}}@{}}
\toprule\noalign{}
\begin{minipage}[b]{\linewidth}\raggedright
target
\end{minipage} & \begin{minipage}[b]{\linewidth}\raggedright
baseline
\end{minipage} & \begin{minipage}[b]{\linewidth}\raggedright
delta\_full
\end{minipage} & \begin{minipage}[b]{\linewidth}\raggedright
ci\_low
\end{minipage} & \begin{minipage}[b]{\linewidth}\raggedright
ci\_high
\end{minipage} & \begin{minipage}[b]{\linewidth}\raggedright
prop\_candidate\_better
\end{minipage} \\
\midrule\noalign{}
\endhead
\bottomrule\noalign{}
\endlastfoot
count & Nd linear & -0.422 & -0.686 & -0.024 & 1.000 \\
count & Nd quadratic & -0.065 & -0.292 & 0.109 & 0.704 \\
count & & Nd-5 & quadratic & 0.004 & -0.018 \\
span & Nd linear & -0.495 & -0.633 & -0.416 & 1.000 \\
span & Nd quadratic & -0.034 & -0.522 & 0.306 & 0.704 \\
span & & Nd-5 & quadratic & -0.001 & -0.062 \\
\end{longtable}

The interpretation is now stable:

\[
\boxed{d_{\mathrm{balance}}\ \text{beats baselines that cannot represent a peak, but not the fair curved filling baselines.}}
\]

For count, the difference against \(N_d+N_d^2\) crosses zero. Against
quadratic \(|N_d-5|\), it is essentially a tie. For span, the same
pattern holds.

\section{4. What survives from the d-block
result}\label{what-survives-from-the-d-block-result}

The d-block experiment should be described as a compact quantitative
restatement of known chemistry:

\[
\boxed{d_{\mathrm{balance}}\ \text{encodes the half-filled-shell symmetry, not a new independent chemical law.}}
\]

The useful parts are:

\begin{itemize}
\tightlist
\item
  the effect is organized by symmetry about \(d^5\);
\item
  it is not well described by strictly monotonic properties alone;
\item
  it provides a clean sanity check for the framework;
\item
  it forces a falsifiable distinction between compact representation and
  predictive novelty.
\end{itemize}

The not-yet-supported parts are:

\begin{itemize}
\tightlist
\item
  ``out-predicts electron count'' as a headline;
\item
  ``new predictor'' relative to quadratic filling;
\item
  any claim that RGM/ROSI is validated by the d-block dataset.
\end{itemize}

\section{5. Associated experiments and current
status}\label{associated-experiments-and-current-status}

\begin{longtable}[]{@{}
  >{\raggedright\arraybackslash}p{(\columnwidth - 8\tabcolsep) * \real{0.2000}}
  >{\raggedright\arraybackslash}p{(\columnwidth - 8\tabcolsep) * \real{0.2000}}
  >{\raggedright\arraybackslash}p{(\columnwidth - 8\tabcolsep) * \real{0.2000}}
  >{\raggedright\arraybackslash}p{(\columnwidth - 8\tabcolsep) * \real{0.2000}}
  >{\raggedright\arraybackslash}p{(\columnwidth - 8\tabcolsep) * \real{0.2000}}@{}}
\toprule\noalign{}
\begin{minipage}[b]{\linewidth}\raggedright
experiment\_id
\end{minipage} & \begin{minipage}[b]{\linewidth}\raggedright
name
\end{minipage} & \begin{minipage}[b]{\linewidth}\raggedright
main\_candidate
\end{minipage} & \begin{minipage}[b]{\linewidth}\raggedright
strongest\_baseline
\end{minipage} & \begin{minipage}[b]{\linewidth}\raggedright
status
\end{minipage} \\
\midrule\noalign{}
\endhead
\bottomrule\noalign{}
\endlastfoot
E1 & Closure-node d-block seed (Zenodo 20573845) & d\_balance = 1- &
Nd-5 & /5 \\
E2 & Chemical-property confound audit & d\_balance added after monotonic
property quadratics & EN+EN\^{}2, IE+IE\^{}2, radius+radius\^{}2 & not a
proxy for monotonic properties; still textbook filling symmetry \\
E3 & Actinide f-balance stress test & f\_balance around f7 & Nf +
Nf\^{}2 & negative/partial stress test \\
E4 & Actinide s-d-f static fraction test & P3, entropy, electron-hole
terms & Nf + Nf\^{}2 / Nf(14-Nf) & static orbital fractions
insufficient \\
E5 & Quadratic radius-complement state/edge screen & R=sqrt(x),
Rbar=sqrt(1-x) & DOF-matched polynomial in same occupancy fractions;
state-output edge baseline & interpretive representation only; predictor
claim falsified in this screen \\
E6 & External redox-couple potential test & radius-complement deltas & Z
+ charge/delta-charge & first non-circular predictive test negative \\
\end{longtable}

\section{6. Post-audit conclusion for radius--complement
geometry}\label{post-audit-conclusion-for-radiuscomplement-geometry}

The radius--complement construction

\[
R_i=\sqrt{x_i},\qquad \bar R_i=\sqrt{1-x_i},\qquad R_i^2+\bar R_i^2=1
\]

is interpretively attractive because it separates occupation from
complement. The product \(R_i\bar R_i\) highlights half filling without
manually adding a half-filled rule.

But the post-audit result is negative:

\[
\boxed{\text{The transformation does not add predictive information beyond flexible polynomials in the same occupancies.}}
\]

The constructed transition-edge task is also not independent when the
label is defined as the conjunction of two endpoint state labels.
Therefore it cannot validate a transition geometry without external
redox-couple data.

\section{7. First non-circular redox-couple
test}\label{first-non-circular-redox-couple-test}

The first non-circular test used redox-couple potentials for U, Np, Pu,
and Am in acidic aqueous solution. The target was no longer ``state
observed'' or ``edge observed'', but an external potential \(E^0\).

The result is negative: simple \(Z\)-and-charge baselines outperform
radius--complement deltas in leave-one-element-out testing.

This closes the current predictive claim:

\[
\boxed{\text{No current experiment establishes predictive novelty for the closure-node or radius--complement descriptors.}}
\]

\section{8. What the article should now
say}\label{what-the-article-should-now-say}

A publication-quality article should state:

\begin{quote}
The closure-node and radius--complement constructions are generative
relational coordinates. They can propose hypotheses, organize known
patterns, and define falsifiable experiments. In the present experiments
they do not yet outperform fair baselines, so they should not be
presented as validated predictors.
\end{quote}

Recommended framing:

\[
\boxed{\text{generation of falsifiable mechanisms, not prediction.}}
\]

\section{9. Response to lyo\_bench}\label{response-to-lyo_bench}

A concise public response can be:

\begin{quote}
You were right on the fair baseline and the bootstrap granularity. I
reran the model comparison with quadratic filling baselines and a
series-block bootstrap. The descriptor still robustly beats linear
\(N_d\), but it does not beat \(N_d+N_d^2\) or quadratic \(|N_d-5|\).
The current status is therefore: compact encoding of half-filled-shell
symmetry, not a predictor established beyond the fair curved filling
baseline. I also agree that the oxidation-state source must be locked
before any stronger claim. The follow-up actinide and radius--complement
tests were treated as falsifier-driven screens; they did not establish
predictive novelty. The framework is useful as a generator of hypotheses
and tests, not yet as a validated chemical predictor.
\end{quote}

\section{10. Final conclusion}\label{final-conclusion}

The most honest closure is:

\[
\boxed{\text{The framework has generated useful coordinates and falsifiers, but the predictor claim is currently negative.}}
\]

The scientific value at this stage is not a successful predictor. It is
the disciplined conversion of a relational intuition into reproducible
tests, fair baselines, and falsifiable next steps.

\section{References}\label{references}

\begin{enumerate}
\def\labelenumi{\arabic{enumi}.}
\tightlist
\item
  L. D. Mata Sánchez, \emph{Closure-Node Registration}, Zenodo record
  \texttt{10.5281/zenodo.20573845}, 2026.
\item
  \texttt{ROSI\_article\_support\_2026-06-06.zip}, supplementary support
  archive attached to Zenodo record \texttt{10.5281/zenodo.20573845}.
\item
  Public lyo\_bench Reddit critique, treated here as an audit prompt
  rather than a bibliographic source.
\end{enumerate}
